% Generated draft: requires mathematical and source review.
\documentclass[11pt]{article}
\usepackage[margin=25mm]{geometry}
\usepackage{fontspec}
\setmainfont{Latin Modern Roman}
\usepackage{amsmath,amssymb,mathtools,unicode-math}
\setmathfont{Latin Modern Math}
\usepackage{enumitem,needspace,xurl,hyperref}
\setlength{\parindent}{0pt}
\setlength{\parskip}{5pt}
\setlength{\emergencystretch}{3em}
\newcommand{\N}{\mathbb{N}}
\newcommand{\Z}{\mathbb{Z}}
\newcommand{\Q}{\mathbb{Q}}
\newcommand{\R}{\mathbb{R}}
\newcommand{\C}{\mathbb{C}}
\newcommand{\F}{\mathbb{F}}
\newcommand{\A}{\mathbb{A}}
\newcommand{\PP}{\mathbb{P}}
\newcommand{\E}{\mathbb{E}}
\newcommand{\eps}{\varepsilon}
\newcommand{\dd}{\,\mathrm{d}}
\begin{document}
\section*{UnsolvedMath ID: 30000109}
\textbf{Source problem code:} catalogue entry
\subsection{Definitions and mathematical statement}
Let $P$ be an integral polytope of affine dimension $m$, using its affine lattice for lattice-point counts. Define the Ehrhart numerator by
\[
 \sum_{r\ge0}\#(rP\cap\Z^m)z^r
   =\frac{h_P^*(z)}{(1-z)^{m+1}},\qquad
 h_P^*(z)=\sum_{i=0}^{s}h_i^*(P)z^i,\quad s\le m,
\]
after choosing full-dimensional lattice coordinates. The full Ehrhart semigroup ring is
\[
 R_P=k[\{(x,r)\in\Z^m\times\Z_{\ge0}:x\in rP\}],
\]
graded by $r$. Standard means generated in degree one; this is a substantive condition on this full semigroup, rather than the automatic standard grading of a ring defined only from degree-one generators. Gorenstein has its usual graded commutative-algebra meaning. A sequence $(a_0,\ldots,a_s)$ is unimodal if some $j$ satisfies $a_0\le\cdots\le a_j\ge\cdots\ge a_s$.

The source asks two separate questions:
\begin{enumerate}
\item If $R_P$ is standard and Gorenstein, must $(h_i^*(P))$ be unimodal?
\item More generally, if the Ehrhart numerator is symmetric, written in the source as $h_i^*(P)=h_{d-i}^*(P)$ for all $i$, must it be unimodal?
\end{enumerate}
For symmetry of the polynomial itself, one takes $d=s=\deg h_P^*$ and interprets the equation for $0\le i\le s$. If instead the source intends the length-$m+1$ vector with trailing zeros and $d=m$, that is a different symmetry convention and must be checked; no equality of degree and dimension is silently assumed.

The literature assessment checked on 21 August 2026 reports counterexamples to the general symmetric-$h^*$ assertion. It separately says that the standard-Gorenstein-semigroup-ring assertion remains open beyond important special classes. The broader counterexamples therefore cannot be presented as counterexamples satisfying the first question's standardness hypothesis without checking them. These are supplied dated status claims, not independently updated literature conclusions.
\subsection{Short English statement}
Must a standard Gorenstein Ehrhart ring have a unimodal $h^*$-vector, and does symmetry alone force unimodality? The supplied assessment distinguishes the first question from the second, which it reports refuted in general.
\subsection{Sources}
\begin{itemize}
\item \emph{Ehrhart Quasipolynomials: Algebra, Combinatorics, and Geometry} (2005), Oberwolfach Reports, DOI 10.4171/owr/2004/39, pp. 2071--2102, \url{https://doi.org/10.4171/owr/2004/39}; proposed year supplied as 2004.
\item Literature assessment checked 21 August 2026: Mustata--Payne, \url{https://web.ma.utexas.edu/users/sampayne/pdf/Ehrhart-polynomials.pdf}; Bruns--Romer, \url{https://arxiv.org/abs/math/0508392}; ``OWR report,'' \url{https://doi.org/10.4171/owr/2004/39}.
\item Local UnsolvedMath registry, permanent ID $30000109$, exact source code \texttt{catalogue entry}, created 25 August 2026 and updated 16 September 2026. Its clean statement is labelled corrected and verified, retaining both assertions and their different hypotheses.
\end{itemize}
\end{document}
